% from main.tex
%%%%%%%%%%  scan idx 135  |  folio 122  %%%%%%%%%%
\emph{avec \(F'\) strictement \(\mathcal C_0\)-parfait (2.1), et tel que \(F'^i=0\) pour \(i\notin[a,b]\), est un raffinement de} \(S\). (\emph{En vertu de (4.3) cette condition équivaut à celle obtenue en remplaçant l'expression « quasi-isomorphisme de \(\Cpx(\mathcal C_X)\) » par « isomorphisme de \(\Dpx(\mathcal C_X)\) ».}) \emph{Nous dirons que \(F\) est d'}amplitude \(\mathcal C_0\)-parfaite finie \emph{s'il existe des entiers \(a\) et \(b\) tels que} \(\mathcal C_0\)-\(\mathrm{parf}\)-\(\amp(F)\subset[a,b]\). \emph{Nous dirons que \(F\) est} \(\mathcal C_0\)-parfait \emph{si \(F\) est localement d'amplitude \(\mathcal C_0\)-parfaite finie}\footnote{Bien entendu, quand cela ne pourra prêter à confusion, nous omettrons \(\mathcal C_0\) dans les expressions précédentes.}.

\textbf{Exemple 4.8.} Soient \((S,A)\) un topos annelé, \(\mathcal C\) la \(S\)-catégorie définie par \(U\longmapsto\) (catégorie des \(A_U\)-Modules à gauche), \(\mathcal C_0\) la sous-\(S\)-catégorie définie par \(U\longmapsto\) (catégorie des \(A_U\)-Modules à gauche facteurs directs de Modules libres de type fini). Le couple \((\mathcal C,\mathcal C_0)\) satisfait, d'après (1.3.1), aux hypothèses (4.0). Un objet \(F\) de \(\Dpx(_A S)\) (cf. (2.15)) sera dit d'amplitude parfaite contenue dans \([a,b]\) (resp. d'amplitude parfaite finie, resp. parfait) si \(F\) est d'amplitude \(\mathcal C_0\)-parfaite contenue dans \([a,b]\) (resp.\dots). Lorsque les facteurs directs de Modules libres de type fini sont localement libres (cf. (2.15.1)), dire que \(F\) est parfait signifie donc que \(F\) est localement isomorphe (dans \(\Dpx(S)\)) à un complexe à degrés bornés et à composantes libres de type fini.

\textbf{Notation 4.9.} Nous désignerons par \(\Dpx(\mathcal C)_{\mathcal C_0\text{-}\mathrm{parf}}\) ou \(\Parf(\mathcal C,\mathcal C_0)\) la sous-\(S\)-catégorie strictement pleine de \(\Dpx(\mathcal C)\) formée des objets \(\mathcal C_0\)-parfaits. Dans le cas de l'exemple (4.8), nous écrirons \(\Dpx(_A S)_{\mathrm{parf}}\), ou \(\Parf(_A S)\), ou même \(\Parf(S)\) s'il n'y a pas d'ambiguïté possible sur \(A\).

